% !TeX program = xelatex
\documentclass[UTF8,a4paper,zihao=-4,fontset=windows,oneside,openany]{ctexbook}
\usepackage[margin=22mm,headheight=15pt]{geometry}
\usepackage{amsmath,amssymb,graphicx,xcolor,needspace,longtable,fancyhdr}
\usepackage[unicode,hidelinks]{hyperref}
\setlength{\parindent}{0pt}
\setlength{\parskip}{4pt}
\setlength{\emergencystretch}{3em}
\linespread{1.2}
\raggedbottom
\pagestyle{fancy}
\fancyhf{}
\fancyhead[L]{湘教版数学必修第一册：教师转写工作稿}
\fancyhead[R]{题干与公式可编辑}
\fancyfoot[C]{\thepage}
\renewcommand{\headrulewidth}{0.3pt}
\newcommand{\subq}[1]{\par\Needspace{2\baselineskip}\noindent（#1）\quad}
\newcommand{\choice}[1]{\par\noindent #1.\quad}
\newcommand{\blank}{\underline{\hspace{1.7cm}}}
\newcommand{\ocrglyph}[1]{\textcolor{red!70!black}{[识别字符 U+#1 待核]}}
\newenvironment{bankitem}[4]{\par\Needspace{5\baselineskip}\noindent
  {\bfseries #1}\par{\small\color{blue!55!black}#2}\par
  {\footnotesize\ttfamily\detokenize{#3}}\par
  {\footnotesize\color{red!65!black}#4}\par\nopagebreak}
  {\par\medskip\hrule\medskip}
\newenvironment{ocrdraft}{\begingroup\small\color{black!65}}{\par\endgroup}
\newcommand{\questioninput}[1]{\input{questions/#1.tex}}
\hypersetup{pdftitle={湘教版必修第一册教材题库：转写工作稿},pdfauthor={教学资料整理}}

\begin{document}
\frontmatter
\begin{titlepage}\centering\vspace*{22mm}
{\Huge 湘教版高中数学\\[5mm]必修第一册教材题库\par}
\vspace{14mm}{\Large 题干与公式可编辑的转写工作稿\par}
\vspace{12mm}{\large 757 个已定位条目；期中范围368个条目已对照原页转写\par}
\vspace{10mm}\begin{minipage}{0.88\textwidth}
\textbf{进度说明：}第1至3章及第4章4.1、4.2（教材至第114页）共368个条目的题干及公式已转写为数学 TeX，配图单独保存，数据表可编辑；其余389个条目仍为 OCR 草稿。原页中的真假命题、判断选项按原文保留，不表示其结论成立。
\par 本工作稿用于教师核查，不是学生卷，不附作答留白，不含独立答案。例题原解及相邻说明保存在旁边的原图题库中，不混入已转写题干。
\par 每个条目独立保存在 questions 目录，可通过 \texttt{\textbackslash questioninput\{题目ID\}} 引用。必需图形单独保存在 figures；未转写条目的图表尚待核对。
\par 757为主标签定位的题目或题组数，小问不拆分。页边提问框、阅读材料内的散见问句尚未逐项列入完整性审查；不据此声称全部教材设问已无遗漏。
\par 印刷页码与 PDF 页序不同：正文 PDF 页序=教材页码+7。源教材未修改；所有条目仍待教师复核后用于课堂。
\end{minipage}\vfill{\large 2026年10月}\end{titlepage}
\tableofcontents\mainmatter

\chapter{集合与逻辑}
\questioninput{XJ-B1-C01-TH-P002-Y0577}
\questioninput{XJ-B1-C01-EX-P003-Y0620}
\questioninput{XJ-B1-C01-EX-P003-Y1223}
\questioninput{XJ-B1-C01-PR-P004-Y0218}
\questioninput{XJ-B1-C01-PR-P004-Y0443}
\questioninput{XJ-B1-C01-EX-P004-Y1108}
\questioninput{XJ-B1-C01-EX-P005-Y0448}
\questioninput{XJ-B1-C01-EX-P005-Y1339}
\questioninput{XJ-B1-C01-PR-P006-Y0217}
\questioninput{XJ-B1-C01-PR-P006-Y0346}
\questioninput{XJ-B1-C01-PR-P006-Y0475}
\questioninput{XJ-B1-C01-TH-P006-Y0849}
\questioninput{XJ-B1-C01-EX-P007-Y0656}
\questioninput{XJ-B1-C01-TH-P007-Y1130}
\questioninput{XJ-B1-C01-EX-P008-Y0280}
\questioninput{XJ-B1-C01-EX-P008-Y0515}
\questioninput{XJ-B1-C01-PR-P008-Y0877}
\questioninput{XJ-B1-C01-PR-P008-Y0920}
\questioninput{XJ-B1-C01-PR-P008-Y1092}
\questioninput{XJ-B1-C01-PR-P008-Y1178}
\questioninput{XJ-B1-C01-TH-P009-Y0951}
\questioninput{XJ-B1-C01-EX-P010-Y0276}
\questioninput{XJ-B1-C01-EX-P010-Y0596}
\questioninput{XJ-B1-C01-TH-P010-Y1233}
\questioninput{XJ-B1-C01-EX-P011-Y0417}
\questioninput{XJ-B1-C01-EX-P011-Y0583}
\questioninput{XJ-B1-C01-TH-P011-Y1054}
\questioninput{XJ-B1-C01-PR-P011-Y1439}
\questioninput{XJ-B1-C01-PR-P011-Y1484}
\questioninput{XJ-B1-C01-PR-P011-Y1526}
\questioninput{XJ-B1-C01-PR-P011-Y1570}
\questioninput{XJ-B1-C01-HW-P012-Y0440}
\questioninput{XJ-B1-C01-HW-P012-Y0651}
\questioninput{XJ-B1-C01-HW-P012-Y0852}
\questioninput{XJ-B1-C01-HW-P012-Y1024}
\questioninput{XJ-B1-C01-HW-P012-Y1154}
\questioninput{XJ-B1-C01-HW-P012-Y1286}
\questioninput{XJ-B1-C01-HW-P012-Y1412}
\questioninput{XJ-B1-C01-HW-P013-Y0307}
\questioninput{XJ-B1-C01-HW-P013-Y0433}
\questioninput{XJ-B1-C01-HW-P013-Y0665}
\questioninput{XJ-B1-C01-HW-P013-Y0750}
\questioninput{XJ-B1-C01-HW-P013-Y0794}
\questioninput{XJ-B1-C01-HW-P013-Y0922}
\questioninput{XJ-B1-C01-HW-P013-Y1051}
\questioninput{XJ-B1-C01-HW-P013-Y1352}
\questioninput{XJ-B1-C01-HW-P013-Y1442}
\questioninput{XJ-B1-C01-HW-P013-Y1531}
\questioninput{XJ-B1-C01-TH-P014-Y1066}
\questioninput{XJ-B1-C01-EX-P015-Y0177}
\questioninput{XJ-B1-C01-EX-P015-Y0738}
\questioninput{XJ-B1-C01-PR-P016-Y0217}
\questioninput{XJ-B1-C01-PR-P016-Y0429}
\questioninput{XJ-B1-C01-PR-P016-Y0642}
\questioninput{XJ-B1-C01-TH-P016-Y0927}
\questioninput{XJ-B1-C01-EX-P017-Y0928}
\questioninput{XJ-B1-C01-EX-P017-Y1338}
\questioninput{XJ-B1-C01-EX-P018-Y0491}
\questioninput{XJ-B1-C01-PR-P019-Y0218}
\questioninput{XJ-B1-C01-PR-P019-Y0389}
\questioninput{XJ-B1-C01-PR-P019-Y0648}
\questioninput{XJ-B1-C01-TH-P019-Y1033}
\questioninput{XJ-B1-C01-EX-P020-Y0562}
\questioninput{XJ-B1-C01-TH-P020-Y0942}
\questioninput{XJ-B1-C01-EX-P020-Y1143}
\questioninput{XJ-B1-C01-EX-P021-Y0177}
\questioninput{XJ-B1-C01-PR-P021-Y0793}
\questioninput{XJ-B1-C01-PR-P021-Y1067}
\questioninput{XJ-B1-C01-TH-P021-Y1462}
\questioninput{XJ-B1-C01-EX-P022-Y0891}
\questioninput{XJ-B1-C01-EX-P022-Y1296}
\questioninput{XJ-B1-C01-PR-P023-Y0176}
\questioninput{XJ-B1-C01-HW-P023-Y0778}
\questioninput{XJ-B1-C01-HW-P023-Y0993}
\questioninput{XJ-B1-C01-HW-P023-Y1165}
\questioninput{XJ-B1-C01-HW-P023-Y1381}
\questioninput{XJ-B1-C01-HW-P024-Y0260}
\questioninput{XJ-B1-C01-HW-P024-Y0518}
\questioninput{XJ-B1-C01-HW-P024-Y0897}
\questioninput{XJ-B1-C01-HW-P024-Y1155}
\questioninput{XJ-B1-C01-HW-P024-Y1370}
\questioninput{XJ-B1-C01-HW-P024-Y1527}
\questioninput{XJ-B1-C01-SU-P028-Y0974}
\questioninput{XJ-B1-C01-SU-P028-Y1103}
\questioninput{XJ-B1-C01-SU-P028-Y1232}
\questioninput{XJ-B1-C01-SU-P028-Y1361}
\questioninput{XJ-B1-C01-SU-P028-Y1404}
\questioninput{XJ-B1-C01-SU-P028-Y1490}
\questioninput{XJ-B1-C01-RV-P029-Y0406}
\questioninput{XJ-B1-C01-RV-P029-Y0634}
\questioninput{XJ-B1-C01-RV-P029-Y0779}
\questioninput{XJ-B1-C01-RV-P029-Y0908}
\questioninput{XJ-B1-C01-RV-P029-Y0994}
\questioninput{XJ-B1-C01-RV-P029-Y1080}
\questioninput{XJ-B1-C01-RV-P029-Y1165}
\questioninput{XJ-B1-C01-RV-P029-Y1294}
\questioninput{XJ-B1-C01-RV-P029-Y1552}
\questioninput{XJ-B1-C01-RV-P030-Y0307}
\questioninput{XJ-B1-C01-RV-P030-Y0393}
\questioninput{XJ-B1-C01-RV-P030-Y0711}
\questioninput{XJ-B1-C01-RV-P030-Y0753}
\questioninput{XJ-B1-C01-RV-P030-Y0836}
\questioninput{XJ-B1-C01-RV-P030-Y0883}
\questioninput{XJ-B1-C01-RV-P030-Y1051}
\questioninput{XJ-B1-C01-RV-P030-Y1286}
\questioninput{XJ-B1-C01-RV-P030-Y1329}
\questioninput{XJ-B1-C01-RV-P030-Y1372}

\chapter{一元二次函数、方程和不等式}
\questioninput{XJ-B1-C02-TH-P032-Y0904}
\questioninput{XJ-B1-C02-CA-P033-Y0234}
\questioninput{XJ-B1-C02-CA-P033-Y0951}
\questioninput{XJ-B1-C02-EX-P034-Y0478}
\questioninput{XJ-B1-C02-EX-P034-Y0712}
\questioninput{XJ-B1-C02-PR-P034-Y1377}
\questioninput{XJ-B1-C02-PR-P034-Y1563}
\questioninput{XJ-B1-C02-TH-P035-Y0319}
\questioninput{XJ-B1-C02-EX-P037-Y0431}
\questioninput{XJ-B1-C02-EX-P037-Y0590}
\questioninput{XJ-B1-C02-PR-P037-Y0875}
\questioninput{XJ-B1-C02-PR-P037-Y1148}
\questioninput{XJ-B1-C02-PR-P037-Y1391}
\questioninput{XJ-B1-C02-EX-P039-Y0303}
\questioninput{XJ-B1-C02-EX-P039-Y0870}
\questioninput{XJ-B1-C02-EX-P039-Y1344}
\questioninput{XJ-B1-C02-PR-P040-Y0217}
\questioninput{XJ-B1-C02-PR-P040-Y0349}
\questioninput{XJ-B1-C02-PR-P040-Y0406}
\questioninput{XJ-B1-C02-PR-P040-Y0461}
\questioninput{XJ-B1-C02-TH-P040-Y0846}
\questioninput{XJ-B1-C02-EX-P041-Y0698}
\questioninput{XJ-B1-C02-EX-P041-Y1378}
\questioninput{XJ-B1-C02-PR-P042-Y0646}
\questioninput{XJ-B1-C02-PR-P042-Y1011}
\questioninput{XJ-B1-C02-PR-P042-Y1470}
\questioninput{XJ-B1-C02-HW-P043-Y0450}
\questioninput{XJ-B1-C02-HW-P043-Y0493}
\questioninput{XJ-B1-C02-HW-P043-Y0634}
\questioninput{XJ-B1-C02-HW-P043-Y0819}
\questioninput{XJ-B1-C02-HW-P043-Y1062}
\questioninput{XJ-B1-C02-HW-P043-Y1406}
\questioninput{XJ-B1-C02-HW-P043-Y1506}
\questioninput{XJ-B1-C02-HW-P044-Y0220}
\questioninput{XJ-B1-C02-HW-P044-Y0732}
\questioninput{XJ-B1-C02-HW-P044-Y0861}
\questioninput{XJ-B1-C02-HW-P044-Y1102}
\questioninput{XJ-B1-C02-HW-P044-Y1420}
\questioninput{XJ-B1-C02-TH-P045-Y0479}
\questioninput{XJ-B1-C02-EX-P046-Y1243}
\questioninput{XJ-B1-C02-EX-P047-Y0578}
\questioninput{XJ-B1-C02-PR-P047-Y1177}
\questioninput{XJ-B1-C02-PR-P047-Y1307}
\questioninput{XJ-B1-C02-PR-P047-Y1435}
\questioninput{XJ-B1-C02-PR-P047-Y1521}
\questioninput{XJ-B1-C02-HW-P049-Y0449}
\questioninput{XJ-B1-C02-HW-P049-Y0579}
\questioninput{XJ-B1-C02-HW-P049-Y0622}
\questioninput{XJ-B1-C02-HW-P049-Y0793}
\questioninput{XJ-B1-C02-HW-P049-Y0879}
\questioninput{XJ-B1-C02-HW-P049-Y1197}
\questioninput{XJ-B1-C02-HW-P049-Y1284}
\questioninput{XJ-B1-C02-TH-P050-Y1033}
\questioninput{XJ-B1-C02-EX-P052-Y0592}
\questioninput{XJ-B1-C02-EX-P052-Y1407}
\questioninput{XJ-B1-C02-EX-P053-Y0304}
\questioninput{XJ-B1-C02-PR-P053-Y1221}
\questioninput{XJ-B1-C02-PR-P053-Y1335}
\questioninput{XJ-B1-C02-EX-P054-Y1398}
\questioninput{XJ-B1-C02-EX-P055-Y0177}
\questioninput{XJ-B1-C02-EX-P055-Y0626}
\questioninput{XJ-B1-C02-PR-P055-Y1064}
\questioninput{XJ-B1-C02-PR-P055-Y1178}
\questioninput{XJ-B1-C02-PR-P055-Y1260}
\questioninput{XJ-B1-C02-PR-P055-Y1345}
\questioninput{XJ-B1-C02-EX-P056-Y0386}
\questioninput{XJ-B1-C02-EX-P056-Y1002}
\questioninput{XJ-B1-C02-EX-P057-Y0534}
\questioninput{XJ-B1-C02-PR-P058-Y0217}
\questioninput{XJ-B1-C02-PR-P058-Y0432}
\questioninput{XJ-B1-C02-HW-P058-Y1137}
\questioninput{XJ-B1-C02-HW-P058-Y1265}
\questioninput{XJ-B1-C02-HW-P058-Y1395}
\questioninput{XJ-B1-C02-HW-P059-Y0221}
\questioninput{XJ-B1-C02-HW-P059-Y0410}
\questioninput{XJ-B1-C02-HW-P059-Y0467}
\questioninput{XJ-B1-C02-HW-P059-Y0788}
\questioninput{XJ-B1-C02-HW-P059-Y0874}
\questioninput{XJ-B1-C02-HW-P059-Y0917}
\questioninput{XJ-B1-C02-HW-P059-Y1115}
\questioninput{XJ-B1-C02-HW-P059-Y1199}
\questioninput{XJ-B1-C02-HW-P059-Y1286}
\questioninput{XJ-B1-C02-SU-P060-Y0997}
\questioninput{XJ-B1-C02-SU-P060-Y1083}
\questioninput{XJ-B1-C02-SU-P060-Y1169}
\questioninput{XJ-B1-C02-SU-P060-Y1255}
\questioninput{XJ-B1-C02-SU-P060-Y1427}
\questioninput{XJ-B1-C02-RV-P061-Y0406}
\questioninput{XJ-B1-C02-RV-P061-Y0650}
\questioninput{XJ-B1-C02-RV-P061-Y0779}
\questioninput{XJ-B1-C02-RV-P061-Y1157}
\questioninput{XJ-B1-C02-RV-P061-Y1244}
\questioninput{XJ-B1-C02-RV-P061-Y1386}
\questioninput{XJ-B1-C02-RV-P061-Y1530}
\questioninput{XJ-B1-C02-RV-P062-Y0430}
\questioninput{XJ-B1-C02-RV-P062-Y0642}
\questioninput{XJ-B1-C02-RV-P062-Y0684}
\questioninput{XJ-B1-C02-RV-P062-Y0894}
\questioninput{XJ-B1-C02-RV-P062-Y0994}
\questioninput{XJ-B1-C02-RV-P062-Y1179}
\questioninput{XJ-B1-C02-RV-P062-Y1398}
\questioninput{XJ-B1-C02-RV-P063-Y0289}
\questioninput{XJ-B1-C02-RV-P063-Y0863}
\questioninput{XJ-B1-C02-RV-P063-Y1035}

\chapter{函数的概念与性质}
\questioninput{XJ-B1-C03-TH-P065-Y0594}
\questioninput{XJ-B1-C03-TH-P065-Y1561}
\questioninput{XJ-B1-C03-EX-P066-Y0934}
\questioninput{XJ-B1-C03-EX-P066-Y1436}
\questioninput{XJ-B1-C03-EX-P067-Y0440}
\questioninput{XJ-B1-C03-PR-P067-Y0874}
\questioninput{XJ-B1-C03-PR-P067-Y1075}
\questioninput{XJ-B1-C03-PR-P067-Y1287}
\questioninput{XJ-B1-C03-PR-P067-Y1404}
\questioninput{XJ-B1-C03-PR-P067-Y1533}
\questioninput{XJ-B1-C03-EX-P069-Y0904}
\questioninput{XJ-B1-C03-EX-P071-Y1368}
\questioninput{XJ-B1-C03-PR-P072-Y1158}
\questioninput{XJ-B1-C03-PR-P072-Y1330}
\questioninput{XJ-B1-C03-PR-P072-Y1416}
\questioninput{XJ-B1-C03-LAB-P073-Y0145}
\questioninput{XJ-B1-C03-EX-P074-Y0394}
\questioninput{XJ-B1-C03-EX-P075-Y0177}
\questioninput{XJ-B1-C03-EX-P075-Y0639}
\questioninput{XJ-B1-C03-PR-P076-Y0217}
\questioninput{XJ-B1-C03-PR-P076-Y0364}
\questioninput{XJ-B1-C03-PR-P076-Y0567}
\questioninput{XJ-B1-C03-HW-P076-Y1136}
\questioninput{XJ-B1-C03-HW-P076-Y1523}
\questioninput{XJ-B1-C03-HW-P077-Y0228}
\questioninput{XJ-B1-C03-HW-P077-Y0286}
\questioninput{XJ-B1-C03-HW-P077-Y0403}
\questioninput{XJ-B1-C03-HW-P077-Y0601}
\questioninput{XJ-B1-C03-HW-P077-Y0687}
\questioninput{XJ-B1-C03-HW-P077-Y0773}
\questioninput{XJ-B1-C03-HW-P077-Y1227}
\questioninput{XJ-B1-C03-HW-P077-Y1356}
\questioninput{XJ-B1-C03-HW-P078-Y0473}
\questioninput{XJ-B1-C03-HW-P078-Y0673}
\questioninput{XJ-B1-C03-HW-P078-Y0717}
\questioninput{XJ-B1-C03-HW-P078-Y0808}
\questioninput{XJ-B1-C03-HW-P078-Y1137}
\questioninput{XJ-B1-C03-HW-P078-Y1352}
\questioninput{XJ-B1-C03-TH-P079-Y0868}
\questioninput{XJ-B1-C03-TH-P081-Y0178}
\questioninput{XJ-B1-C03-EX-P081-Y0375}
\questioninput{XJ-B1-C03-EX-P081-Y1313}
\questioninput{XJ-B1-C03-EX-P082-Y0618}
\questioninput{XJ-B1-C03-PR-P082-Y1064}
\questioninput{XJ-B1-C03-PR-P082-Y1193}
\questioninput{XJ-B1-C03-PR-P082-Y1536}
\questioninput{XJ-B1-C03-TH-P083-Y0277}
\questioninput{XJ-B1-C03-TH-P083-Y0867}
\questioninput{XJ-B1-C03-EX-P084-Y0177}
\questioninput{XJ-B1-C03-EX-P084-Y0592}
\questioninput{XJ-B1-C03-PR-P084-Y1009}
\questioninput{XJ-B1-C03-PR-P084-Y1402}
\questioninput{XJ-B1-C03-PR-P084-Y1488}
\questioninput{XJ-B1-C03-PR-P084-Y1532}
\questioninput{XJ-B1-C03-HW-P085-Y0449}
\questioninput{XJ-B1-C03-HW-P085-Y0836}
\questioninput{XJ-B1-C03-HW-P085-Y1036}
\questioninput{XJ-B1-C03-HW-P085-Y1122}
\questioninput{XJ-B1-C03-HW-P085-Y1280}
\questioninput{XJ-B1-C03-HW-P086-Y0231}
\questioninput{XJ-B1-C03-HW-P086-Y0288}
\questioninput{XJ-B1-C03-HW-P086-Y0520}
\questioninput{XJ-B1-C03-HW-P086-Y0708}
\questioninput{XJ-B1-C03-HW-P086-Y0853}
\questioninput{XJ-B1-C03-HW-P086-Y1034}
\questioninput{XJ-B1-C03-HW-P086-Y1121}
\questioninput{XJ-B1-C03-HW-P086-Y1208}
\questioninput{XJ-B1-C03-SU-P089-Y0854}
\questioninput{XJ-B1-C03-SU-P089-Y1026}
\questioninput{XJ-B1-C03-SU-P089-Y1112}
\questioninput{XJ-B1-C03-SU-P089-Y1284}
\questioninput{XJ-B1-C03-SU-P089-Y1370}
\questioninput{XJ-B1-C03-RV-P090-Y0407}
\questioninput{XJ-B1-C03-RV-P090-Y0635}
\questioninput{XJ-B1-C03-RV-P090-Y0678}
\questioninput{XJ-B1-C03-RV-P090-Y0720}
\questioninput{XJ-B1-C03-RV-P090-Y1103}
\questioninput{XJ-B1-C03-RV-P090-Y1174}
\questioninput{XJ-B1-C03-RV-P090-Y1256}
\questioninput{XJ-B1-C03-RV-P090-Y1299}
\questioninput{XJ-B1-C03-RV-P090-Y1401}
\questioninput{XJ-B1-C03-RV-P090-Y1501}
\questioninput{XJ-B1-C03-RV-P090-Y1543}
\questioninput{XJ-B1-C03-RV-P091-Y0650}
\questioninput{XJ-B1-C03-RV-P091-Y0785}
\questioninput{XJ-B1-C03-RV-P091-Y1267}
\questioninput{XJ-B1-C03-RV-P091-Y1481}
\questioninput{XJ-B1-C03-RV-P092-Y0231}
\questioninput{XJ-B1-C03-RV-P092-Y0288}
\questioninput{XJ-B1-C03-RV-P092-Y0375}
\questioninput{XJ-B1-C03-RV-P092-Y1067}
\questioninput{XJ-B1-C03-RV-P092-Y1491}

\chapter{幂函数、指数函数和对数函数}
\questioninput{XJ-B1-C04-EX-P095-Y0782}
\questioninput{XJ-B1-C04-EX-P096-Y0533}
\questioninput{XJ-B1-C04-EX-P096-Y0978}
\questioninput{XJ-B1-C04-EX-P096-Y1306}
\questioninput{XJ-B1-C04-EX-P097-Y0403}
\questioninput{XJ-B1-C04-PR-P097-Y0891}
\questioninput{XJ-B1-C04-PR-P097-Y1020}
\questioninput{XJ-B1-C04-PR-P097-Y1240}
\questioninput{XJ-B1-C04-PR-P097-Y1468}
\questioninput{XJ-B1-C04-TH-P098-Y0419}
\questioninput{XJ-B1-C04-TH-P098-Y1358}
\questioninput{XJ-B1-C04-EX-P099-Y1332}
\questioninput{XJ-B1-C04-EX-P099-Y1526}
\questioninput{XJ-B1-C04-PR-P100-Y0459}
\questioninput{XJ-B1-C04-PR-P100-Y0574}
\questioninput{XJ-B1-C04-PR-P100-Y0660}
\questioninput{XJ-B1-C04-EX-P102-Y0177}
\questioninput{XJ-B1-C04-EX-P102-Y0669}
\questioninput{XJ-B1-C04-PR-P103-Y0797}
\questioninput{XJ-B1-C04-PR-P103-Y1012}
\questioninput{XJ-B1-C04-PR-P103-Y1370}
\questioninput{XJ-B1-C04-PR-P103-Y1471}
\questioninput{XJ-B1-C04-HW-P104-Y0450}
\questioninput{XJ-B1-C04-HW-P104-Y0598}
\questioninput{XJ-B1-C04-HW-P104-Y0788}
\questioninput{XJ-B1-C04-HW-P104-Y0874}
\questioninput{XJ-B1-C04-HW-P104-Y1201}
\questioninput{XJ-B1-C04-HW-P104-Y1259}
\questioninput{XJ-B1-C04-HW-P104-Y1352}
\questioninput{XJ-B1-C04-HW-P104-Y1494}
\questioninput{XJ-B1-C04-HW-P105-Y0348}
\questioninput{XJ-B1-C04-HW-P105-Y0405}
\questioninput{XJ-B1-C04-HW-P105-Y0490}
\questioninput{XJ-B1-C04-HW-P105-Y0578}
\questioninput{XJ-B1-C04-HW-P105-Y0620}
\questioninput{XJ-B1-C04-TH-P106-Y0720}
\questioninput{XJ-B1-C04-EX-P108-Y0177}
\questioninput{XJ-B1-C04-EX-P108-Y0840}
\questioninput{XJ-B1-C04-PR-P108-Y1273}
\questioninput{XJ-B1-C04-PR-P108-Y1402}
\questioninput{XJ-B1-C04-TH-P109-Y0302}
\questioninput{XJ-B1-C04-EX-P110-Y0319}
\questioninput{XJ-B1-C04-TH-P110-Y1359}
\questioninput{XJ-B1-C04-EX-P111-Y1422}
\questioninput{XJ-B1-C04-EX-P112-Y0737}
\questioninput{XJ-B1-C04-EX-P112-Y1015}
\questioninput{XJ-B1-C04-PR-P113-Y0217}
\questioninput{XJ-B1-C04-PR-P113-Y0303}
\questioninput{XJ-B1-C04-PR-P113-Y0347}
\questioninput{XJ-B1-C04-PR-P113-Y0475}
\questioninput{XJ-B1-C04-HW-P113-Y0886}
\questioninput{XJ-B1-C04-HW-P113-Y1057}
\questioninput{XJ-B1-C04-HW-P113-Y1144}
\questioninput{XJ-B1-C04-HW-P113-Y1315}
\questioninput{XJ-B1-C04-HW-P113-Y1446}
\questioninput{XJ-B1-C04-HW-P113-Y1566}
\questioninput{XJ-B1-C04-HW-P114-Y0303}
\questioninput{XJ-B1-C04-HW-P114-Y0667}
\questioninput{XJ-B1-C04-HW-P114-Y0711}
\questioninput{XJ-B1-C04-HW-P114-Y0935}
\questioninput{XJ-B1-C04-HW-P114-Y1038}
\questioninput{XJ-B1-C04-HW-P114-Y1222}
\questioninput{XJ-B1-C04-HW-P114-Y1332}
\questioninput{XJ-B1-C04-HW-P114-Y1396}
\questioninput{XJ-B1-C04-HW-P114-Y1507}
\questioninput{XJ-B1-C04-TH-P115-Y0591}
\questioninput{XJ-B1-C04-EX-P115-Y1132}
\questioninput{XJ-B1-C04-EX-P116-Y0405}
\questioninput{XJ-B1-C04-PR-P116-Y0981}
\questioninput{XJ-B1-C04-PR-P116-Y1252}
\questioninput{XJ-B1-C04-PR-P116-Y1455}
\questioninput{XJ-B1-C04-EX-P117-Y1349}
\questioninput{XJ-B1-C04-EX-P118-Y0345}
\questioninput{XJ-B1-C04-EX-P118-Y0571}
\questioninput{XJ-B1-C04-PR-P118-Y1113}
\questioninput{XJ-B1-C04-PR-P118-Y1234}
\questioninput{XJ-B1-C04-PR-P118-Y1421}
\questioninput{XJ-B1-C04-TH-P119-Y0670}
\questioninput{XJ-B1-C04-EX-P119-Y1195}
\questioninput{XJ-B1-C04-EX-P119-Y1481}
\questioninput{XJ-B1-C04-EX-P120-Y0421}
\questioninput{XJ-B1-C04-EX-P120-Y0749}
\questioninput{XJ-B1-C04-PR-P121-Y0218}
\questioninput{XJ-B1-C04-PR-P121-Y0261}
\questioninput{XJ-B1-C04-PR-P121-Y0304}
\questioninput{XJ-B1-C04-PR-P121-Y0344}
\questioninput{XJ-B1-C04-EX-P122-Y1083}
\questioninput{XJ-B1-C04-EX-P122-Y1421}
\questioninput{XJ-B1-C04-EX-P123-Y0664}
\questioninput{XJ-B1-C04-PR-P124-Y0217}
\questioninput{XJ-B1-C04-PR-P124-Y0304}
\questioninput{XJ-B1-C04-PR-P124-Y0490}
\questioninput{XJ-B1-C04-PR-P124-Y0633}
\questioninput{XJ-B1-C04-HW-P124-Y1051}
\questioninput{XJ-B1-C04-HW-P124-Y1208}
\questioninput{XJ-B1-C04-HW-P124-Y1336}
\questioninput{XJ-B1-C04-HW-P124-Y1493}
\questioninput{XJ-B1-C04-HW-P125-Y0416}
\questioninput{XJ-B1-C04-HW-P125-Y0488}
\questioninput{XJ-B1-C04-HW-P125-Y0545}
\questioninput{XJ-B1-C04-HW-P125-Y0804}
\questioninput{XJ-B1-C04-HW-P125-Y0933}
\questioninput{XJ-B1-C04-HW-P125-Y1091}
\questioninput{XJ-B1-C04-HW-P125-Y1228}
\questioninput{XJ-B1-C04-HW-P125-Y1398}
\questioninput{XJ-B1-C04-HW-P126-Y0336}
\questioninput{XJ-B1-C04-HW-P126-Y0393}
\questioninput{XJ-B1-C04-HW-P126-Y0450}
\questioninput{XJ-B1-C04-HW-P126-Y0534}
\questioninput{XJ-B1-C04-HW-P126-Y0578}
\questioninput{XJ-B1-C04-HW-P126-Y0711}
\questioninput{XJ-B1-C04-HW-P126-Y0974}
\questioninput{XJ-B1-C04-EX-P130-Y1358}
\questioninput{XJ-B1-C04-EX-P131-Y0177}
\questioninput{XJ-B1-C04-EX-P131-Y1058}
\questioninput{XJ-B1-C04-PR-P132-Y0219}
\questioninput{XJ-B1-C04-PR-P132-Y0350}
\questioninput{XJ-B1-C04-TH-P132-Y0715}
\questioninput{XJ-B1-C04-EX-P132-Y1265}
\questioninput{XJ-B1-C04-EX-P134-Y0942}
\questioninput{XJ-B1-C04-PR-P135-Y0713}
\questioninput{XJ-B1-C04-PR-P135-Y0800}
\questioninput{XJ-B1-C04-HW-P136-Y0450}
\questioninput{XJ-B1-C04-HW-P136-Y0563}
\questioninput{XJ-B1-C04-HW-P136-Y0620}
\questioninput{XJ-B1-C04-HW-P136-Y0720}
\questioninput{XJ-B1-C04-HW-P136-Y0953}
\questioninput{XJ-B1-C04-HW-P136-Y1124}
\questioninput{XJ-B1-C04-HW-P136-Y1210}
\questioninput{XJ-B1-C04-LAB-P137-Y0145}
\questioninput{XJ-B1-C04-EX-P139-Y0872}
\questioninput{XJ-B1-C04-EX-P140-Y1262}
\questioninput{XJ-B1-C04-PR-P143-Y0217}
\questioninput{XJ-B1-C04-PR-P143-Y0308}
\questioninput{XJ-B1-C04-PR-P143-Y0353}
\questioninput{XJ-B1-C04-EX-P143-Y1169}
\questioninput{XJ-B1-C04-PR-P146-Y0993}
\questioninput{XJ-B1-C04-PR-P146-Y1390}
\questioninput{XJ-B1-C04-HW-P147-Y0449}
\questioninput{XJ-B1-C04-HW-P147-Y0649}
\questioninput{XJ-B1-C04-HW-P147-Y1062}
\questioninput{XJ-B1-C04-HW-P148-Y0218}
\questioninput{XJ-B1-C04-HW-P148-Y0816}
\questioninput{XJ-B1-C04-HW-P148-Y1232}
\questioninput{XJ-B1-C04-HW-P149-Y0347}
\questioninput{XJ-B1-C04-HW-P149-Y0447}
\questioninput{XJ-B1-C04-SU-P150-Y0796}
\questioninput{XJ-B1-C04-SU-P150-Y0882}
\questioninput{XJ-B1-C04-SU-P150-Y1054}
\questioninput{XJ-B1-C04-SU-P150-Y1140}
\questioninput{XJ-B1-C04-SU-P150-Y1226}
\questioninput{XJ-B1-C04-SU-P150-Y1355}
\questioninput{XJ-B1-C04-SU-P150-Y1441}
\questioninput{XJ-B1-C04-RV-P151-Y0407}
\questioninput{XJ-B1-C04-RV-P151-Y0535}
\questioninput{XJ-B1-C04-RV-P151-Y0883}
\questioninput{XJ-B1-C04-RV-P151-Y0949}
\questioninput{XJ-B1-C04-RV-P151-Y0995}
\questioninput{XJ-B1-C04-RV-P151-Y1201}
\questioninput{XJ-B1-C04-RV-P151-Y1258}
\questioninput{XJ-B1-C04-RV-P151-Y1345}
\questioninput{XJ-B1-C04-RV-P151-Y1430}
\questioninput{XJ-B1-C04-RV-P152-Y0219}
\questioninput{XJ-B1-C04-RV-P152-Y0261}
\questioninput{XJ-B1-C04-RV-P152-Y0465}
\questioninput{XJ-B1-C04-RV-P152-Y0508}
\questioninput{XJ-B1-C04-RV-P152-Y0552}
\questioninput{XJ-B1-C04-RV-P152-Y0638}
\questioninput{XJ-B1-C04-RV-P152-Y0721}
\questioninput{XJ-B1-C04-RV-P152-Y0954}
\questioninput{XJ-B1-C04-RV-P152-Y1473}
\questioninput{XJ-B1-C04-RV-P153-Y0321}
\questioninput{XJ-B1-C04-RV-P153-Y0465}
\questioninput{XJ-B1-C04-RV-P153-Y0579}
\questioninput{XJ-B1-C04-RV-P153-Y0875}
\questioninput{XJ-B1-C04-RV-P153-Y1090}
\questioninput{XJ-B1-C04-RV-P153-Y1172}
\questioninput{XJ-B1-C04-RV-P153-Y1216}

\chapter{三角函数}
\questioninput{XJ-B1-C05-TH-P156-Y0591}
\questioninput{XJ-B1-C05-TH-P157-Y0550}
\questioninput{XJ-B1-C05-EX-P157-Y0962}
\questioninput{XJ-B1-C05-EX-P157-Y1411}
\questioninput{XJ-B1-C05-PR-P158-Y0217}
\questioninput{XJ-B1-C05-PR-P158-Y0260}
\questioninput{XJ-B1-C05-PR-P158-Y0346}
\questioninput{XJ-B1-C05-PR-P158-Y0432}
\questioninput{XJ-B1-C05-TH-P158-Y1003}
\questioninput{XJ-B1-C05-TH-P159-Y1004}
\questioninput{XJ-B1-C05-EX-P159-Y1414}
\questioninput{XJ-B1-C05-EX-P160-Y0177}
\questioninput{XJ-B1-C05-EX-P160-Y1053}
\questioninput{XJ-B1-C05-PR-P161-Y0221}
\questioninput{XJ-B1-C05-PR-P161-Y0388}
\questioninput{XJ-B1-C05-PR-P161-Y0517}
\questioninput{XJ-B1-C05-HW-P161-Y1048}
\questioninput{XJ-B1-C05-HW-P161-Y1135}
\questioninput{XJ-B1-C05-HW-P161-Y1177}
\questioninput{XJ-B1-C05-HW-P161-Y1334}
\questioninput{XJ-B1-C05-HW-P161-Y1378}
\questioninput{XJ-B1-C05-HW-P161-Y1464}
\questioninput{XJ-B1-C05-HW-P162-Y0217}
\questioninput{XJ-B1-C05-HW-P162-Y0656}
\questioninput{XJ-B1-C05-HW-P162-Y0743}
\questioninput{XJ-B1-C05-HW-P162-Y1029}
\questioninput{XJ-B1-C05-HW-P162-Y1089}
\questioninput{XJ-B1-C05-HW-P162-Y1218}
\questioninput{XJ-B1-C05-HW-P162-Y1346}
\questioninput{XJ-B1-C05-TH-P163-Y0689}
\questioninput{XJ-B1-C05-EX-P164-Y0486}
\questioninput{XJ-B1-C05-EX-P164-Y0865}
\questioninput{XJ-B1-C05-PR-P165-Y0810}
\questioninput{XJ-B1-C05-PR-P165-Y0853}
\questioninput{XJ-B1-C05-TH-P165-Y1147}
\questioninput{XJ-B1-C05-EX-P166-Y0573}
\questioninput{XJ-B1-C05-TH-P166-Y0848}
\questioninput{XJ-B1-C05-EX-P167-Y0891}
\questioninput{XJ-B1-C05-PR-P167-Y1326}
\questioninput{XJ-B1-C05-PR-P167-Y1442}
\questioninput{XJ-B1-C05-PR-P167-Y1555}
\questioninput{XJ-B1-C05-TH-P168-Y0290}
\questioninput{XJ-B1-C05-EX-P168-Y1068}
\questioninput{XJ-B1-C05-EX-P169-Y0177}
\questioninput{XJ-B1-C05-EX-P169-Y0600}
\questioninput{XJ-B1-C05-EX-P169-Y0756}
\questioninput{XJ-B1-C05-PR-P169-Y1245}
\questioninput{XJ-B1-C05-PR-P169-Y1316}
\questioninput{XJ-B1-C05-PR-P169-Y1372}
\questioninput{XJ-B1-C05-PR-P169-Y1430}
\questioninput{XJ-B1-C05-PR-P169-Y1502}
\questioninput{XJ-B1-C05-EX-P170-Y0796}
\questioninput{XJ-B1-C05-TH-P170-Y0981}
\questioninput{XJ-B1-C05-EX-P172-Y0384}
\questioninput{XJ-B1-C05-EX-P172-Y0880}
\questioninput{XJ-B1-C05-PR-P172-Y1270}
\questioninput{XJ-B1-C05-PR-P172-Y1457}
\questioninput{XJ-B1-C05-PR-P172-Y1512}
\questioninput{XJ-B1-C05-TH-P173-Y0178}
\questioninput{XJ-B1-C05-EX-P174-Y0694}
\questioninput{XJ-B1-C05-EX-P174-Y0981}
\questioninput{XJ-B1-C05-PR-P174-Y1188}
\questioninput{XJ-B1-C05-PR-P174-Y1316}
\questioninput{XJ-B1-C05-PR-P174-Y1444}
\questioninput{XJ-B1-C05-HW-P175-Y0444}
\questioninput{XJ-B1-C05-HW-P175-Y0530}
\questioninput{XJ-B1-C05-HW-P175-Y0644}
\questioninput{XJ-B1-C05-HW-P175-Y0829}
\questioninput{XJ-B1-C05-HW-P175-Y0958}
\questioninput{XJ-B1-C05-HW-P175-Y1016}
\questioninput{XJ-B1-C05-HW-P175-Y1129}
\questioninput{XJ-B1-C05-HW-P175-Y1318}
\questioninput{XJ-B1-C05-HW-P175-Y1379}
\questioninput{XJ-B1-C05-HW-P175-Y1506}
\questioninput{XJ-B1-C05-HW-P176-Y0350}
\questioninput{XJ-B1-C05-HW-P176-Y0407}
\questioninput{XJ-B1-C05-HW-P176-Y0450}
\questioninput{XJ-B1-C05-HW-P176-Y0577}
\questioninput{XJ-B1-C05-HW-P176-Y0704}
\questioninput{XJ-B1-C05-TH-P177-Y0904}
\questioninput{XJ-B1-C05-TH-P178-Y1217}
\questioninput{XJ-B1-C05-EX-P179-Y0233}
\questioninput{XJ-B1-C05-PR-P180-Y0179}
\questioninput{XJ-B1-C05-TH-P180-Y0494}
\questioninput{XJ-B1-C05-EX-P182-Y0945}
\questioninput{XJ-B1-C05-EX-P182-Y1522}
\questioninput{XJ-B1-C05-EX-P183-Y0458}
\questioninput{XJ-B1-C05-PR-P183-Y1029}
\questioninput{XJ-B1-C05-PR-P183-Y1185}
\questioninput{XJ-B1-C05-PR-P183-Y1299}
\questioninput{XJ-B1-C05-PR-P183-Y1484}
\questioninput{XJ-B1-C05-TH-P184-Y1253}
\questioninput{XJ-B1-C05-EX-P185-Y0596}
\questioninput{XJ-B1-C05-EX-P185-Y1160}
\questioninput{XJ-B1-C05-PR-P186-Y0217}
\questioninput{XJ-B1-C05-PR-P186-Y0260}
\questioninput{XJ-B1-C05-PR-P186-Y0375}
\questioninput{XJ-B1-C05-HW-P186-Y0851}
\questioninput{XJ-B1-C05-HW-P186-Y0964}
\questioninput{XJ-B1-C05-HW-P186-Y1094}
\questioninput{XJ-B1-C05-HW-P186-Y1207}
\questioninput{XJ-B1-C05-HW-P186-Y1392}
\questioninput{XJ-B1-C05-HW-P186-Y1522}
\questioninput{XJ-B1-C05-HW-P187-Y0217}
\questioninput{XJ-B1-C05-HW-P187-Y0260}
\questioninput{XJ-B1-C05-HW-P187-Y0452}
\questioninput{XJ-B1-C05-HW-P187-Y0625}
\questioninput{XJ-B1-C05-HW-P187-Y0896}
\questioninput{XJ-B1-C05-HW-P187-Y0986}
\questioninput{XJ-B1-C05-HW-P187-Y1229}
\questioninput{XJ-B1-C05-HW-P187-Y1329}
\questioninput{XJ-B1-C05-TH-P189-Y0178}
\questioninput{XJ-B1-C05-EX-P189-Y0403}
\questioninput{XJ-B1-C05-EX-P190-Y0361}
\questioninput{XJ-B1-C05-EX-P190-Y1515}
\questioninput{XJ-B1-C05-PR-P191-Y1332}
\questioninput{XJ-B1-C05-PR-P191-Y1518}
\questioninput{XJ-B1-C05-EX-P192-Y0318}
\questioninput{XJ-B1-C05-TH-P192-Y1420}
\questioninput{XJ-B1-C05-EX-P193-Y0513}
\questioninput{XJ-B1-C05-EX-P194-Y1243}
\questioninput{XJ-B1-C05-PR-P195-Y0633}
\questioninput{XJ-B1-C05-PR-P195-Y0861}
\questioninput{XJ-B1-C05-PR-P195-Y1290}
\questioninput{XJ-B1-C05-HW-P196-Y1132}
\questioninput{XJ-B1-C05-HW-P196-Y1319}
\questioninput{XJ-B1-C05-HW-P197-Y0418}
\questioninput{XJ-B1-C05-HW-P197-Y0661}
\questioninput{XJ-B1-C05-HW-P197-Y0889}
\questioninput{XJ-B1-C05-HW-P197-Y1160}
\questioninput{XJ-B1-C05-HW-P197-Y1298}
\questioninput{XJ-B1-C05-HW-P197-Y1405}
\questioninput{XJ-B1-C05-LAB-P198-Y0145}
\questioninput{XJ-B1-C05-EX-P199-Y0577}
\questioninput{XJ-B1-C05-EX-P200-Y0177}
\questioninput{XJ-B1-C05-PR-P201-Y0218}
\questioninput{XJ-B1-C05-PR-P201-Y0602}
\questioninput{XJ-B1-C05-HW-P201-Y1132}
\questioninput{XJ-B1-C05-HW-P201-Y1548}
\questioninput{XJ-B1-C05-HW-P202-Y0419}
\questioninput{XJ-B1-C05-HW-P202-Y0822}
\questioninput{XJ-B1-C05-HW-P202-Y0993}
\questioninput{XJ-B1-C05-SU-P205-Y0907}
\questioninput{XJ-B1-C05-SU-P205-Y0993}
\questioninput{XJ-B1-C05-SU-P205-Y1036}
\questioninput{XJ-B1-C05-SU-P205-Y1165}
\questioninput{XJ-B1-C05-SU-P205-Y1294}
\questioninput{XJ-B1-C05-SU-P205-Y1380}
\questioninput{XJ-B1-C05-RV-P206-Y0406}
\questioninput{XJ-B1-C05-RV-P206-Y0564}
\questioninput{XJ-B1-C05-RV-P206-Y0736}
\questioninput{XJ-B1-C05-RV-P206-Y0821}
\questioninput{XJ-B1-C05-RV-P206-Y0923}
\questioninput{XJ-B1-C05-RV-P206-Y0979}
\questioninput{XJ-B1-C05-RV-P206-Y1107}
\questioninput{XJ-B1-C05-RV-P206-Y1435}
\questioninput{XJ-B1-C05-RV-P206-Y1539}
\questioninput{XJ-B1-C05-RV-P207-Y0321}
\questioninput{XJ-B1-C05-RV-P207-Y0378}
\questioninput{XJ-B1-C05-RV-P207-Y0563}
\questioninput{XJ-B1-C05-RV-P207-Y0788}
\questioninput{XJ-B1-C05-RV-P207-Y0960}
\questioninput{XJ-B1-C05-RV-P207-Y1129}
\questioninput{XJ-B1-C05-RV-P208-Y0354}
\questioninput{XJ-B1-C05-RV-P208-Y0426}
\questioninput{XJ-B1-C05-RV-P208-Y0569}
\questioninput{XJ-B1-C05-RV-P208-Y0667}
\questioninput{XJ-B1-C05-RV-P208-Y0802}
\questioninput{XJ-B1-C05-RV-P208-Y0927}
\questioninput{XJ-B1-C05-RV-P208-Y1276}
\questioninput{XJ-B1-C05-RV-P209-Y0333}
\questioninput{XJ-B1-C05-RV-P209-Y1208}

\chapter{统计学初步}
\questioninput{XJ-B1-C06-TH-P211-Y0813}
\questioninput{XJ-B1-C06-PR-P213-Y1211}
\questioninput{XJ-B1-C06-PR-P213-Y1382}
\questioninput{XJ-B1-C06-HW-P214-Y0453}
\questioninput{XJ-B1-C06-HW-P214-Y0496}
\questioninput{XJ-B1-C06-HW-P214-Y0667}
\questioninput{XJ-B1-C06-HW-P214-Y0836}
\questioninput{XJ-B1-C06-HW-P214-Y1026}
\questioninput{XJ-B1-C06-TH-P215-Y1331}
\questioninput{XJ-B1-C06-TH-P217-Y1023}
\questioninput{XJ-B1-C06-PR-P217-Y1376}
\questioninput{XJ-B1-C06-PR-P217-Y1462}
\questioninput{XJ-B1-C06-TH-P218-Y0289}
\questioninput{XJ-B1-C06-EX-P218-Y1196}
\questioninput{XJ-B1-C06-EX-P219-Y1028}
\questioninput{XJ-B1-C06-PR-P220-Y0217}
\questioninput{XJ-B1-C06-PR-P220-Y0475}
\questioninput{XJ-B1-C06-HW-P221-Y1223}
\questioninput{XJ-B1-C06-HW-P221-Y1309}
\questioninput{XJ-B1-C06-HW-P221-Y1395}
\questioninput{XJ-B1-C06-HW-P221-Y1567}
\questioninput{XJ-B1-C06-HW-P222-Y0303}
\questioninput{XJ-B1-C06-HW-P222-Y0664}
\questioninput{XJ-B1-C06-HW-P222-Y0836}
\questioninput{XJ-B1-C06-HW-P222-Y1183}
\questioninput{XJ-B1-C06-CA-P224-Y0720}
\questioninput{XJ-B1-C06-CA-P226-Y0277}
\questioninput{XJ-B1-C06-PR-P227-Y0217}
\questioninput{XJ-B1-C06-PR-P227-Y0307}
\questioninput{XJ-B1-C06-CA-P227-Y1094}
\questioninput{XJ-B1-C06-PR-P230-Y0946}
\questioninput{XJ-B1-C06-PR-P230-Y0988}
\questioninput{XJ-B1-C06-PR-P230-Y1031}
\questioninput{XJ-B1-C06-HW-P231-Y0446}
\questioninput{XJ-B1-C06-HW-P231-Y0872}
\questioninput{XJ-B1-C06-HW-P231-Y1224}
\questioninput{XJ-B1-C06-HW-P232-Y0260}
\questioninput{XJ-B1-C06-HW-P232-Y0932}
\questioninput{XJ-B1-C06-HW-P233-Y0261}
\questioninput{XJ-B1-C06-HW-P233-Y0871}
\questioninput{XJ-B1-C06-HW-P233-Y1402}
\questioninput{XJ-B1-C06-LAB-P234-Y0145}
\questioninput{XJ-B1-C06-EX-P236-Y0767}
\questioninput{XJ-B1-C06-EX-P236-Y1549}
\questioninput{XJ-B1-C06-EX-P237-Y1511}
\questioninput{XJ-B1-C06-PR-P238-Y1159}
\questioninput{XJ-B1-C06-PR-P238-Y1480}
\questioninput{XJ-B1-C06-EX-P239-Y0957}
\questioninput{XJ-B1-C06-EX-P240-Y0489}
\questioninput{XJ-B1-C06-PR-P241-Y0647}
\questioninput{XJ-B1-C06-PR-P241-Y1232}
\questioninput{XJ-B1-C06-TH-P242-Y0889}
\questioninput{XJ-B1-C06-EX-P243-Y1308}
\questioninput{XJ-B1-C06-EX-P244-Y0613}
\questioninput{XJ-B1-C06-EX-P246-Y0361}
\questioninput{XJ-B1-C06-PR-P247-Y0217}
\questioninput{XJ-B1-C06-PR-P247-Y0570}
\questioninput{XJ-B1-C06-PR-P247-Y0922}
\questioninput{XJ-B1-C06-TH-P248-Y0430}
\questioninput{XJ-B1-C06-EX-P248-Y0969}
\questioninput{XJ-B1-C06-TH-P249-Y0892}
\questioninput{XJ-B1-C06-PR-P250-Y0217}
\questioninput{XJ-B1-C06-PR-P250-Y0799}
\questioninput{XJ-B1-C06-TH-P251-Y0637}
\questioninput{XJ-B1-C06-EX-P251-Y1103}
\questioninput{XJ-B1-C06-EX-P252-Y0348}
\questioninput{XJ-B1-C06-PR-P253-Y0668}
\questioninput{XJ-B1-C06-PR-P253-Y0840}
\questioninput{XJ-B1-C06-HW-P254-Y0444}
\questioninput{XJ-B1-C06-HW-P254-Y0953}
\questioninput{XJ-B1-C06-HW-P254-Y1370}
\questioninput{XJ-B1-C06-HW-P255-Y0260}
\questioninput{XJ-B1-C06-HW-P255-Y0519}
\questioninput{XJ-B1-C06-HW-P255-Y0772}
\questioninput{XJ-B1-C06-HW-P255-Y1208}
\questioninput{XJ-B1-C06-HW-P255-Y1508}
\questioninput{XJ-B1-C06-HW-P256-Y0260}
\questioninput{XJ-B1-C06-HW-P256-Y0708}
\questioninput{XJ-B1-C06-HW-P256-Y1125}
\questioninput{XJ-B1-C06-HW-P257-Y0217}
\questioninput{XJ-B1-C06-HW-P257-Y0449}
\questioninput{XJ-B1-C06-HW-P257-Y0885}
\questioninput{XJ-B1-C06-SU-P262-Y1160}
\questioninput{XJ-B1-C06-SU-P262-Y1246}
\questioninput{XJ-B1-C06-SU-P262-Y1504}
\questioninput{XJ-B1-C06-SU-P263-Y0217}
\questioninput{XJ-B1-C06-SU-P263-Y0561}
\questioninput{XJ-B1-C06-RV-P264-Y0406}
\questioninput{XJ-B1-C06-RV-P264-Y0707}
\questioninput{XJ-B1-C06-RV-P264-Y0965}
\questioninput{XJ-B1-C06-RV-P264-Y1051}
\questioninput{XJ-B1-C06-RV-P264-Y1266}
\questioninput{XJ-B1-C06-RV-P265-Y0217}
\questioninput{XJ-B1-C06-RV-P265-Y0501}
\questioninput{XJ-B1-C06-RV-P265-Y0964}
\questioninput{XJ-B1-C06-RV-P265-Y1455}
\questioninput{XJ-B1-C06-RV-P266-Y0326}
\questioninput{XJ-B1-C06-RV-P266-Y0893}
\questioninput{XJ-B1-C06-RV-P266-Y1140}
\questioninput{XJ-B1-C06-RV-P267-Y0346}
\questioninput{XJ-B1-C06-RV-P267-Y1120}
\questioninput{XJ-B1-C06-RV-P268-Y0432}
\questioninput{XJ-B1-C06-RV-P268-Y1202}
\questioninput{XJ-B1-C06-RV-P269-Y0453}

\end{document}
